\documentclass[11pt]{article}

% --- Core Academic Packages ---
\usepackage[utf8]{inputenc}
\usepackage{amsmath, amsthm, amssymb, amsfonts} % Essential for formal econometric math
\usepackage{booktabs}                           % Publication-quality tables (\toprule, \midrule, \bottomrule)
\usepackage{threeparttable}                    % Crucial for robust standard error notes below tables
\usepackage{graphicx}                          % For visual proofs/plots
\usepackage{subcaption}                        % Panel A / Panel B structural layouts
\usepackage{geometry}                           % Standard margins
\usepackage{setspace}                           % Double/1.5 spacing for referee drafts
\usepackage[toc,page]{appendix}                % Formal appendix formatting
\usepackage[colorlinks=true, linkcolor=blue, citecolor=blue]{hyperref}
\usepackage[backend=biber, style=authoryear, maxcitenames=2, uniquename=false]{biblatex}
\usepackage{enumitem}

\geometry{a4paper, margin=0.70in}

\addbibresource{bib.bib}

\title{\textbf{Economic Development, a Political Materialism perspective}}
\author{\textbf{Carlos Gustavo Salas Flores} \\ \small Beijing Normal University}
\date{Oct 2026}

\begin{document}

\maketitle


% \newpage

% =========================================================================
% 1. BACKGROUND
% =========================================================================
\section{Background}
\label{sec:background}


Recent Economic models attempt to capture latent political variables such as democracy metrics. Some of the most widely used metrics in academia today are V-Dem and Polity-V; however, \textcite{treier_jackman_2008} showed significant error in Polity-V's democracy measures that lead to significant inference errors. On another study, \textcite{little_meng_2024} challenged recent claims that show democracy backsliding in recent years and found that such claims are not supported by objective data such as election results, competitiveness of elections, and journalists jailed and murdered; such findings suggest that it's the experts' biases that are changing instead of objective measures.


One of the most influential works in recent years has been that of AJR, who developed a theory of Inclusive and Extractive Political Institutions that attempts to resolve why two territories with similar material conditions can develop very differently, as in the case of Nogales, Arizona, and Nogales, Coahuila \parencite{acemoglu_johnson_robinson_2001,acemoglu_johnson_robinson_2005,acemoglu_robinson_2012}. The main drawback of this theory is its inability to develop a theory of the state; instead, it relies on economic indicators and experts’ opinions to classify a nation into Inclusive (II) or Extractive (EI). There are, however, at least two flaws to this model. \textit{I) Fallacy of the converse.} Saying that a country with Inclusive Institutions must carry out a Low Expropriation Rate ($A \to X$) means that, if a country carries out a Low Expropriation Rate, then that country must have Inclusive Institutions ($X \to A$); we cannot assume reverse causality. \textit{II) Vague Horizon.} AJR attempts to resolve the last point by saying that some Extractive Institutions can carry out Inclusive Economic policies, but such economies are destined to fail in the long term. This is an Unfalsifiable Claim or a Vague Horizon; basically, it cannot be proven right or wrong.

To explain the phenomenon described in \textcite{little_meng_2024} and offer an alternative to AJR's theory, I draw from concepts of classic Political and Legal theory. Firstly, we need to make a distinction between the regime and the establishment. A regime refers to the actors who govern (e.g., presidents, Representatives, and Ruling Parties) and any of their actions (Policies and International Relations). The establishment refers to every institution that is mostly funded by the state, such as the Army, Police Force, Public Hospitals, Research Institutes, Schools, and State Companies.

The Leviathan, or the state, is the government and its institutions that maintain order in a given
territory through its legal monopoly of violence \parencite{hobbes1651leviathan,weber1922wirtschaft}. Hence, regime change, as explained in \textcite{acemoglu_naidu_restrepo_robinson_2019}, would result in economic growth. However, such a study relies on democracy indicators rather than directly analyzing state structures. Furthermore, the findings in \textcite{little_meng_2024} cast some doubt on whether these correlations are due to objective changes in such economies or rather changes in the biases of experts' opinions about these economies. An alternative explanation is that regime changes happen without any structural change \parencite{garciatrevijano2010}; hence, the economic and military situation of an economy is rather driven by International Relations \parencite{waltz_1979, mearsheimer_2001}, and conditioned on its material conditions \parencite{riedel_sachs_2005}.

In synthesis, this work attempts to prove the following hypothesis.

\begin{itemize}
    \item $H_1$ on Internal Regime Changes: Any two countries, $A$ and $B$, with the same system of government and similar geographical conditions, will not suffer major economic shocks driven by internal regime changes when under the same sphere of influence by great powers.
    \item $H_2$ on International Politics: Any two countries, $A$ and $B$, with the same system of government and similar geographical conditions, regardless of the regime, will be primarily affected by international relations.
    \item $H_3$ on Simultaneous Regime and IR Changes: This is a logical consequence of $H_1$ and $H_2$. Any country $A$, that preserves its system of government and goes through a regime change will suffer a major economic shock if and only if it also alters its place in the International Political structure by getting in/out a sphere of influence of a major power.
\end{itemize}





% =========================================================================
% 2. LITERATURE REVIEW
% =========================================================================
\section{Literature Review and Theoretical Framework}
\label{sec:lit_review}


% =========================================================================
% 8. CONCLUSIONS & LIMITATIONS
% =========================================================================
\section{Limitations and Conclusion}
\label{sec:discussion}



\printbibliography

\end{document}
